\section{Discussion}

\subsection{Residual accuracy is not solution accuracy}

The most important observation is the separation between a small residual and a
small forward error. For the near-singular diagonal system, the component
associated with $10^{-14}$ contributes only $10^{-14}$ to the unnormalized
residual when it is unresolved. With a tolerance of $10^{-12}$, the solver
therefore returns convergence after one iteration even though
$\eta_x=0.7071$. This is expected conditioning behavior, not a contradiction of
the stopping rule.

The complex Hilbert experiment is less extreme but qualitatively similar:
$\rho=5.610\times10^{-15}$ accompanies a $5.47\%$ relative solution error. A
residual is a backward-error-like observable; it is not the usual normwise
backward error $\lVert b-Ax\rVert_2/(\lVert A\rVert_2\lVert x\rVert_2+
\lVert b\rVert_2)$. For nonsingular $A$, nonzero $x_\star$, and $b=Ax_\star$,
\begin{equation}
    \eta_x\leq\kappa_2(A)\rho(x),\qquad
    \kappa_2(A)=\lVert A\rVert_2\lVert A^{-1}\rVert_2.
\end{equation}
For the tests' rounded right-hand side, write $\delta b=b-Ax_\star$; then
\begin{equation}
    \lVert x-x_\star\rVert_2\leq\lVert A^{-1}\rVert_2
       \bigl(\lVert b-Ax\rVert_2+\lVert\delta b\rVert_2\bigr).
\end{equation}
Both the residual and right-hand-side formation error can therefore be amplified
by conditioning. General finite-precision and conditioning principles are
reviewed by Higham \cite{higham_2002}. Consequently,
future experiment reports must retain both $\rho$ and $\eta_x$ whenever a trusted
reference solution is available.

\subsection{Estimated and true residuals}

The small least-squares problem can yield a much more optimistic residual estimate
than a fresh evaluation of $b-Ax$. For the nonnormal Jordan-like system, the last
estimate was $6.079\times10^{-33}$ while the independently accumulated true
residual was $1.368\times10^{-16}$. In the unscaled $3\times3$ system, the
corresponding pair was $4.284\times10^{-32}$ and
$1.608\times10^{-16}$.

The latter system was repeated with a deliberately stringent tolerance
$10^{-30}$. The estimate crossed the tolerance, but the true residual did not;
the solver returned \texttt{BREAKDOWN} after three iterations rather than
\texttt{CONVERGED}. This tolerance was not attained in the reported run; it is
not claimed to be unattainable for every double-precision system. The maximum
absolute gap between the estimate and the solver's double-precision residual,
$\max_t|\widehat{\rho}_t-\rho_{\mathrm{solver},t}|$, was
$1.665\times10^{-16}$. The independent wide residual is the final value quoted
above, not a history recorded at each step. Thus, in this reproducible example,
the independent residual guard prevents false convergence. It does not eliminate
roundoff from the true-residual computation itself, which remains a double-
precision CSR matrix--vector product.

\subsection{Restart and breakdown behavior}

Restart length changes the available minimization space. The skew-symmetric
$2\times2$ example demonstrates exact GMRES(1) stagnation, whereas restart two
recovers the exact solution in two iterations. Other baseline systems converge
with restart one or two, but require 27 and 19 total iterations, respectively.
These cases verify cycle continuation and global iteration accounting without
supporting a general convergence-rate claim.

The implementation distinguishes a true-residual success at near-breakdown from
breakdown without convergence. Identity-like cases exercise happy breakdown; the
inconsistent zero-operator case returns \texttt{BREAKDOWN} with residual one.
This status communicates that the Krylov process cannot supply another useful
direction under the implemented threshold.
